\subsection{Method}
\label{sec:method}

We study a broader class of \textbf{post-generation selective-deferral selectors}
that operate after candidate generation is complete (Figure~\ref{fig:method_overview}).
Such selectors do not generate new candidates and do not issue additional
inference calls: they choose a final answer from already-produced outputs under
the same per-method budget cap. The scientific object is therefore a routing
problem over answer sources, not a new candidate-generation method.

Failure-Trace Allocator (FTA) is the deterministic instantiation of that broader
framework studied in this paper. Its two gates---the \textit{weak-frontier trigger}
and the \textit{direct-agreement trigger}---fire in strict priority order; at
most one gate acts per example. All runtime decisions are gold-free: no
\texttt{gold\_answer}, \texttt{exact\_match}, or correctness label is accessed.

Throughout the paper we use short manuscript names for readability.
The combined Failure-Trace Allocator (FTA) policy is the \emph{failure-trace-guided allocator}.
Its two internal conditions are the \emph{weak-frontier trigger} and the \emph{direct-agreement trigger}.
The three external baselines are abbreviated \emph{L1}, \emph{S1}, and \emph{TALE}; the best-performing external on a given evaluation split is the \emph{best external}.
Implementation-name mapping is retained in Table~\ref{tab:name_mapping} for reproducibility.

At each example, the selector consumes three classes of runtime-legal inputs: the
frontier answer candidate, the three evaluated external-baseline answers, and
failure-trace metadata emitted during normal frontier execution. The output is a
single final answer plus an explicit policy outcome
(\texttt{fix2}, \texttt{fix4}, or frontier default), so every decision can be
audited after the run using only logged runtime fields.
This interface is intentionally narrow: the instantiated policy does not depend on
learned scores or post-hoc labels.

Conceptually, the allocator uses three abstract predicates.
\textsc{WeakFrontier}$(x)$ is true when the frontier trace indicates low support
or a single-branch selection.
\textsc{ExternalMajority}$(x)$ returns an external answer when at least two of
L1, S1, and TALE agree, with a deterministic fallback only when all three differ.
\textsc{ExternalConsensusDisagree}$(x)$ is true when L1, S1, and TALE unanimously
agree with one another and disagree with the frontier answer.
The raw log strings in Table~\ref{alg:fix24} and Table~\ref{tab:policy_rules}
are implementation realizations of these predicates, not scientific assumptions
tied to a private script.

The weak-frontier trigger is motivated by a structural reliability pattern. In practice, branches and candidate groups are formed from the frontier's logged branch traces: each produced branch yields a parsed answer and is assigned to an answer group by normalized predicted text. Support counts are the integer counts of branches per answer-group (\texttt{answer\_group\_support\_counts}) observed in the trace; higher support indicates multiple independent branches converging on the same answer. The field \texttt{override\_reason} records the runtime decision pathway (for example, \texttt{single\_weak\_frontier\_branch} or \texttt{direct\_frontier\_agree}) and is emitted by the frontier when a guard or post-processing event signals limited search breadth or internal agreement. Weak-frontier traces are emitted when expansion depth or family diversity falls below configured thresholds; direct-frontier agreement is emitted when the frontier's reserve selection matches the incumbent. These metadata fields are runtime-visible and used only for gold-free routing predicates; they are not derived from gold labels or post-hoc scoring.
When the frontier trace indicates a single weak branch (for example, low-depth
termination in the logged metadata), the resulting candidate is treated as
evidence of under-exploration.
In those cases, routing to an external majority answer provides an independent
cross-policy signal without spending extra compute.

The direct-agreement trigger addresses a complementary pattern of selective
deferral.
If the frontier and its reserve agree internally, but all three external methods
unanimously select a different answer, this unanimous disagreement is treated as
high-confidence evidence that the frontier selection may be wrong under the
matched-budget protocol.
Partial external agreement (for example, only one or two externals disagreeing)
is intentionally not used in the main method, because it is more ambiguous in
this setting and did not provide validation-grade evidence.

\medskip\noindent\textbf{Frontier metadata interface (runtime, gold-free).}
The allocator consumes only fields produced during standard frontier execution:
\begin{itemize}[nosep,leftmargin=1.5em]
  \item \textbf{Candidate answer grouping:} parsed branch answers are canonicalized into answer groups (normalization/canonicalization details are implementation-level and documented in reproducibility artifacts).
  \item \textbf{Support counts:} each answer group has an integer branch-support count; these counts define whether support is concentrated or dispersed.
  \item \textbf{Logged pathway metadata:} \texttt{override\_reason} records the frontier selector pathway, including \texttt{single\_weak\_frontier\_branch} and \texttt{direct\_frontier\_agree}.
  \item \textbf{Agreement flags:} direct-frontier agreement and external answer agreement are computed from already generated answers.
\end{itemize}
All of these are runtime-visible model-output artifacts; none require ground-truth answers.

\medskip\noindent\textbf{Conditional-deferral proposition.}
Let $S$ be a runtime-identifiable subset of examples. Let $e_F(S)$ be frontier error rate on $S$ and $e_M(S)$ be external-majority error rate on $S$. If $e_F(S) > e_M(S)$, then replacing frontier with external majority on $S$ improves expected accuracy by
\[
P(S)\,\big(e_F(S)-e_M(S)\big).
\]
Low-Depth Guard (LDG) and External-Consensus Fallback (ECO) are empirical,
gold-free routing predicates intended to identify such subsets without using
correctness labels at runtime.

The control flow is a strict precedence chain (Figure~\ref{fig:method_overview},
Table~\ref{tab:policy_rules}).
The weak-frontier trigger is checked first; when it fires, the policy routes to an
external majority answer and stops.
This priority reflects the interpretation of the weak-frontier trace as a
generation-quality warning: when exploration itself appears narrow, the policy
does not also ask whether the frontier was internally self-consistent.
If it does not fire, the policy checks the direct-agreement trigger; when that
condition fires, the policy routes to the external unanimity answer.
If neither trigger fires, the frontier answer is kept.
The default therefore preserves the original frontier method on all examples
where neither high-confidence failure trace is present.
Conflict handling is deterministic: at most one gate can act, and external
majority ties are resolved by the fixed priority order documented in
Table~\ref{tab:policy_rules}.
That fallback is used for reproducibility rather than claimed as optimal;
the sensitivity noted in Table~\ref{tab:param_sensitivity} shows that alternative
tie orders change Aggregate-720 replay by less than one percentage point.

All trigger features are runtime-legal and gold-free by construction.
They are computed only from model outputs and frontier trace metadata that are
already visible at inference time, never from \texttt{gold\_answer},
\texttt{exact\_match}, example identifiers, or artifact paths.
Table~\ref{tab:feature_legality} gives the complete legality inventory, including
allowed inputs and explicitly excluded fields.
This legality boundary is central to the method definition, not only a reporting
constraint.

Relative to L1, S1, TALE, and the frontier baseline, Failure-Trace Allocator (FTA) plays a different
role in the stack.
L1, S1, TALE, and frontier each provide candidate answers under the same
per-method matched budget; the selector then decides whether to keep the
frontier answer or defer to another answer source using trace and agreement signals.
This separation between candidate generation and final allocation is what allows the
policy to improve selection quality without additional model calls.
Observed gains come from switching the final answer on a small trace-identified
subset of examples rather than from increasing inference budget.
Table~\ref{tab:name_mapping} provides the manuscript-to-implementation name mapping.

\subsection{Low-Depth Guard (LDG)}
\label{sec:fix2}

The Low-Depth Guard detects frontier answers arising from structurally weak search branches and defers to an external majority vote. The trigger relies exclusively on fields logged during the frontier run: the override reason string and the candidate-pool group count.

\subsection{External-Consensus Fallback (ECO)}
\label{sec:fix4}

The External-Consensus Fallback identifies examples where the direct-agreement trigger fires---meaning the frontier agreed with its own internal reserve selection---yet all three external baselines unanimously chose a different answer. When this unanimous disagreement holds, the policy defers to the external consensus answer.

\subsection{Combined Policy Summary}
\label{sec:fix24_algorithm}

Table~\ref{alg:fix24} gives the complete pseudocode using exact runtime field names. Parameter settings and implementation-name mapping are documented in the appendix reproducibility tables.

\begin{table}[t]
\centering
\caption{Policy summary for the Failure-Trace Allocator. All inputs are gold-free runtime fields. Named constants and name mapping are documented in the appendix reproducibility tables.}
\label{alg:fix24}
\scriptsize
\begin{tabular}{p{0.06\textwidth} p{0.88\textwidth}}
\toprule
\textbf{Step} & \textbf{Operation} \\
\midrule
\multicolumn{2}{l}{\textit{Inputs per example}} \\
& \texttt{frontier\_answer}, \texttt{result\_metadata} (from frontier run logs) \\
& \texttt{l1\_answer}, \texttt{s1\_answer}, \texttt{tale\_answer} (external baseline outputs) \\
& \texttt{override\_reason}, \texttt{frontier\_support} $\in \mathbb{Z}_{\geq 0}$, \texttt{candidate\_pool\_answer\_group\_count} $\in \mathbb{Z}_{>0}$ \\
\midrule
\multicolumn{2}{l}{\textit{Step 1 ---Low-Depth Guard (LDG)}} \\
1a & \textbf{if} \texttt{override\_reason} contains \texttt{"single\_weak\_frontier\_branch"} \\
   & \quad \textbf{or} (\texttt{frontier\_support} $= 0$ \textbf{and} \texttt{candidate\_pool\_answer\_group\_count} $\geq 2$) \\
   & \textbf{then:} \\
1b & \quad Compute \texttt{ext} $\leftarrow$ majority answer among \{\texttt{l1}, \texttt{s1}, \texttt{tale}\} (at least 2 agree); \\
   & \quad if no majority, use priority order TALE $\succ$ S1 $\succ$ L1 \\
1c & \quad \textbf{if} \texttt{ext} exists \textbf{and} \texttt{ext} $\neq$ \texttt{frontier\_answer}: \\
   & \quad\quad \textbf{return} \texttt{ext}, \texttt{policy} $\leftarrow$ \texttt{fix2} \quad\quad \textit{(guard fires; stop)} \\
\midrule
\multicolumn{2}{l}{\textit{Step 2 ---External-Consensus Fallback (ECO)}} \\
2a & \textbf{if} \texttt{override\_reason} $=$ \texttt{"direct\_frontier\_agree"} \\
   & \quad \textbf{and} \texttt{l1\_answer} $=$ \texttt{s1\_answer} $=$ \texttt{tale\_answer} \\
   & \quad \textbf{and} \texttt{l1\_answer} $\neq$ \texttt{frontier\_answer} \\
   & \textbf{then:} \\
2b & \quad \textbf{return} \texttt{l1\_answer}, \texttt{policy} $\leftarrow$ \texttt{fix4} \quad\quad \textit{(consensus fallback fires; stop)} \\
\midrule
\multicolumn{2}{l}{\textit{Step 3 ---Default}} \\
3  & \textbf{return} \texttt{frontier\_answer}, \texttt{policy} $\leftarrow$ \texttt{original} \\
\midrule
\multicolumn{2}{l}{\textit{Output}} \\
& Selected answer and \texttt{policy} label; no gold labels accessed at any step. \\
\bottomrule
\end{tabular}
\end{table}

The exact field-level trigger conditions are also tabulated in Table~\ref{tab:policy_rules} for reproducibility.
